\section{Conclusion}

\indent The goal of this project was to understand how different cooling devices could reduce passenger discomfort in outdoor and semi-outdoor waiting areas. More specifically, we wanted to find out how much cooling reaches the waiting area, how that cooling feels to a passenger, and which option performs best across different hot-weather conditions. We compared a fan and two evaporative coolers with a no-device baseline at 12 venues.

\indent In Model 1, we found that the cooling at the outlet of an evaporative cooler is greater than the cooling in the waiting area after the air mixes. The amount of cooling available also depends on outdoor humidity. Model 2 used Standard Effective Temperature (SET) to compare how the devices feel to a passenger and how much discomfort each device reduces per kilowatt-hour. In Model 3, the larger evaporative cooler, Option C, emerged as the preferred choice under the study's comfort objective. Even with Option C, peak-hour SET ranged from $28.12$ to $33.98\,^\circ\mathrm{C}$ and stayed above the $26\,^\circ\mathrm{C}$ comparison target at all 12 venues. Sensitivity analysis found five intervals where a Boston recommendation switched between the fan and the larger cooler. In three of those intervals, increasing the outdoor temperature in the hot-humid scenario favored the fan.

\indent The models connect through the difference between cooling the air and cooling the passenger. The air temperature and humidity calculated in Model 1 become inputs to Model 2, where air movement also matters. This gives the fan a way to reduce discomfort without changing the air temperature. Model 3 then uses the remaining discomfort to compare each option across weather scenarios, and the sensitivity analysis shows where changes in those conditions can change the choice.

\indent For transit agencies, Option C is a starting point where reducing hot discomfort is the main goal. The Boston results give agencies a reason to compare it closely with the fan at humid sites, especially where electricity and water use matter. A useful next step would be to measure airspeed and radiant temperature at passenger positions and check how well the model represents the places where people actually wait. Since peak-hour discomfort remained even with the larger cooler, future work could also test whether pairing a device with better shade brings passengers closer to the comfort target.